% PNP RESEARCH REPORT::ADAPTIVE-COLLISION-CLOSURE-BARRIERS::2026-11-11T20:17:00-08:00
\documentclass[11pt]{article}
\usepackage[T1]{fontenc}
\usepackage[utf8]{inputenc}
\usepackage[margin=1in]{geometry}
\usepackage{amsmath,amssymb,amsthm}
\usepackage{booktabs,longtable,array,xurl}
\usepackage[hidelinks]{hyperref}
\setlength{\emergencystretch}{3em}
\newtheorem{theorem}{Theorem}
\newtheorem{lemma}{Lemma}
\newtheorem{proposition}{Proposition}
\newtheorem{corollary}{Corollary}
\newcommand{\poly}{\operatorname{poly}}
\newcommand{\cl}{\operatorname{cl}}
\newcommand{\F}{\mathbb F_2}
\newcommand{\SAT}{\operatorname{SAT}}
\title{Adaptive Collision Chains and Closure Algorithms:\\
Two Quantitative Obstructions to SAT Work Amplification}
\author{P vs NP Research Project\\Codex research with independent temporary-agent review}
\date{11 November 2026, 20:17 PST}
\begin{document}
\maketitle

\section*{Status and accomplished step}
\textbf{P versus NP remains open.}
This continuation tests the previous report's explicit amplification target
\cite{transfer2026}. The adaptive collision upper bound, the nonlinear
inverse-based construction, the closure lemma, and the stated clock obstruction
are \textbf{PROVED/EXACT} under their explicit computational assumptions.
The finite Boolean and CNF audits are \textbf{COMPUTATIONAL}.
The required lower bound on minimum verified witness work remains \textbf{OPEN}.

The concrete nonlinear candidate fails: truncating a uniformly invertible
keyed permutation permits a collision from two inverse evaluations at every
stage, including when the key depends on the preceding witnesses.
The positive progress is a quantitative exclusion test for later candidates:
small output width and small closure degree both force upper bounds below
the proposed squared-logarithm lower bound. No mathematical priority claim
is made for pigeonhole enumeration, Feistel inversion, or fixed-rank
Horn/affine tractability.

\section*{Model, exact target, and SAME\_REALIZATION}
All logarithms have base two. Inputs are explicitly encoded canonical
3-CNFs with densely indexed occurring variables. Set $N=|F|+2$.
Use the previous report's fixed efficiently universal deterministic machine,
verification allowance $q(N)\ge N+1$, and description budget
$b(N)=b_0+\lfloor\log N\rfloor$:
\[
 H(F)=\min_{|p|\le b(N)}\{t_p(F)+q(N)+1\}.
 \tag{1}
\]
The time is infinite unless $p$ halts with an independently verified
satisfying assignment. A fixed exhaustive finder is admitted.
Constant-size programs eventually fit the budget.

For every fixed correct total SAT decider $A$, the recovered bridge gives
\[
 \overline H(N)\le C_A\bigl(N+T_A(KN)+2\bigr)^D,
 \tag{2}
\]
where $T_A$ is its global worst-case envelope and $K,D$ are fixed
machine/encoding constants. Conversely, an explicit uniform envelope
$\overline H(N)\le h(N)$ yields SAT time $\poly(N+h(N))$ by clocking
the $O(N)$ admitted programs. Hence polynomial boundedness and each fixed
quasipolynomial logarithmic degree transfer in both directions
\cite{transfer2026}. An upper bound for one generated family is not an
upper bound for all SAT.

The target under audit is a satisfiable sequence $F_j$ with
$s_j=|F_j|+2\to\infty$ and fixed $C>1,c>0,B\ge1$ such that
\[
 s_{j+1}\le Cs_j,\qquad
 H(F_{j+1})\ge H(F_j)s_j^c/B.
 \tag{3}
\]
The previous report proves that this would force
\[
 \log H(F_j)\ge
 \frac{c}{2\log C}(\log s_j)^2-O(\log s_j),
 \tag{4}
\]
and a corresponding lower bound on every SAT decider through (2).
Inequality (3) has not been proved. A fixed work multiplier alone does
not yield (4).

All upper bounds below refer to one uniform algorithm constructing the
witness for the exact final formula. Circuit acquisition, prefix
reconstruction, representation extraction, auxiliary gate assignments,
and verification are charged. They are not claims that optimal programs
for separate instances can be concatenated for free under budget (1).

\section{Adaptive total-extension chains have additive search cost}

Let the $i$th stage have an explicit Boolean circuit
\[
 D_i:\{0,1\}^{\ell_i}\times\{0,1\}^{r_i+1}\longrightarrow
       \{0,1\}^{r_i}.
\]
Its first argument $z$ contains the preceding collision pairs.
A stage witness is a pair $u_i\ne v_i$ with
$D_i(z,u_i)=D_i(z,v_i)$. Let $F_k$ encode all stages, their interfaces,
and circuit evaluations using exact Tseitin constraints and standard
proper-3CNF completion.

\begin{theorem}[Adaptive collision replay]\label{thm:collision}
Assume the stage descriptions are explicitly supplied or recoverable by
one uniform polynomial-time procedure from the canonical input family.
Assume the number of stages, circuit sizes, interface lengths, reconstruction
cost, and witness-lifting cost are bounded by one fixed polynomial in the
final length $N$. Assume \emph{every valid prefix} admits the next shrinking
map as specified. Put $R=\max_{i\le k}r_i$. Then a deterministic algorithm
finds a witness in native time
\[
 \poly(N)\sum_{i=1}^k 2^{r_i},
 \tag{5}
\]
and the fixed-machine benchmark satisfies
\[
 H(F_k)\le \poly(N)2^{O(R)}.
 \tag{6}
\]
\end{theorem}

\begin{proof}
At stage $i$, fix the already obtained prefix. Evaluate the circuit on
$2^{r_i}+1$ distinct inputs. There are only $2^{r_i}$ possible outputs,
so sorting the output/input pairs exposes a repeated output with distinct
inputs. Evaluation, sorting, and bookkeeping cost
$\poly(N)2^{r_i}$. Retain one collision and proceed.
Because every chosen valid prefix extends, no earlier collision needs
to be reconsidered. Summing gives (5). Evaluate all circuit gates and
fill the proper-3CNF auxiliary variables to obtain a full satisfying
assignment. One fixed program implements the procedure. The fixed
polynomial universal-simulation overhead and verification allowance
give (6).
\end{proof}

Consequently $R=O(\log N)$ gives polynomial witness work.
If $R=o((\log N)^2)$, (6) is
$2^{o((\log N)^2)}$, contradicting (4).
Within this architecture a squared-logarithm amplifier therefore requires
\[
 R=\Omega((\log N)^2).
 \tag{7}
\]
This is a necessary condition, not sufficient hardness.

The same proof applies to any sequential total search interface whose
next valid extension can be found in time $Q_i(N)$ for \emph{every
selected valid prefix}: the native cost is a polynomial overhead plus
$\sum_i Q_i(N)$. The stronger condition is essential.
Satisfiability preservation alone, meaning that \emph{some} old witness
extends, does not permit this greedy replay. A construction with
incompatible prefixes requires a separately charged compatibility search.

Hubacek's NestedCollision definition is a precise literature counterpart:
find a collision of $f$, then a collision of $g$ after fixing the first
pair. Its adaptive reduction uses two collision calls.
The paper's separation is for shallow decision-tree reductions accessing
function tables through queries \cite{hubacek2026}.
It does not supply an unrestricted time lower bound for explicit
Boolean circuits, and does not turn (5) into a product.

\section{An explicit nonlinear construction that stays easy}

Let $P_z$ be a permutation of $\{0,1\}^{r+1}$ whose forward and inverse
maps are uniformly computable in polynomial time from the prefix-derived
key $z$. Define
\[
 C_z(u)=\operatorname{prefix}_r(P_z(u)).
\]
Then
\[
 u=P_z^{-1}(0^r0),\qquad v=P_z^{-1}(0^r1)
 \tag{8}
\]
are distinct and satisfy $C_z(u)=C_z(v)=0^r$.
This applies at every stage, irrespective of earlier keys.

The implemented control uses four Feistel rounds on two $h$-bit halves.
For each round,
\[
 (L,R)\longmapsto(R,L\oplus f_z(R)),\qquad
 (L',R')\longmapsto(R'\oplus f_z(L'),L').
 \tag{9}
\]
The round function contains an AND of two right-half coordinates,
XORs, prefix-key bits, and round constants. It is nonlinear for the
tested four-bit controls; inversion follows directly from (9), not from
guessing or a cryptographic assumption. The collision map discards one
of the $2h$ output bits, so $r=2h-1$.

The code fixes a three-bit witnessed seed, derives every next key from
the preceding collision inputs, and compiles all XOR/AND/OR gates,
truncated-output equalities, and input inequalities into proper 3-CNF.
Short clauses are expanded with two shared fresh padding variables,
with exact truth-table checks.
For fixed four-round depth, the compiler has
\[
 \begin{aligned}
 \text{gates}&=k(45h-1),\\
 \text{variables}&=7+k(49h-1),\\
 \text{clauses}&=15+k(198h-5).
 \end{aligned}
 \tag{10}
\]
Binary variable identifiers contribute the usual logarithmic encoding
factor. The family is canonical and polynomially reconstructible:
the integer parameters can be enumerated within polynomial bounds
and the generated encoding compared with the input.

Two inverse evaluations per layer construct all primary witnesses;
forward evaluations then validate them and fill gate variables.
Thus the whole family has polynomial $H$, including when output widths
grow beyond (7). High nonlinearity and candidate-dependent sequential keys
do not prevent this shortcut. The quantity recorded as inverse
round-function calls counts the collision-construction phase only,
not all forward checks and CNF completion.

All well-formed instances of the shrinking collision search problem
are satisfiable by construction. Its recognized promise decision
question is therefore trivial. A general SAT decider would still yield
a search algorithm by self-reduction, which queries formulas outside
that promise family. Thus a genuine $H$ lower bound on this family
would remain meaningful, but the present construction has the explicit
upper bound (8).

\section{A closure-degree algorithm for compatible extensions}

The current project note AFFINE-TABLE-COMPOSITION2 proves a closure
enumeration algorithm for Horn constraints plus affine equations
\cite{atc2026}. The following statement exposes the more general
algebraic reason and its quantitative consequence for the amplifier.

\begin{lemma}[Small seeds in an intersection-closed family]\label{lem:seed}
Let $\mathcal M\subseteq 2^{[m]}$ be a nonempty finite family closed under
intersection. For a seed $S\subseteq[m]$, let $\cl(S)$ be the intersection
of all members containing $S$, or $\bot$ if none exists.
Suppose $a:\mathcal M\to\{0,1\}$ is represented on $\mathcal M$ by a
multilinear polynomial over $\F$ of degree at most $d$.
If some $M\in\mathcal M$ has $a(M)=1$, then
\[
 a(\cl(S))=1
 \quad\text{for some }S,\quad |S|\le \min(d,m).
 \tag{11}
\]
\end{lemma}

\begin{proof}
Write the representing polynomial as
$p(x)=\sum_{|S|\le d}c_S\prod_{i\in S}x_i$.
On members $M$ of the family, a monomial is one exactly when
$S\subseteq M$, equivalently $\cl(S)\subseteq M$.
A seed with no extension contributes zero everywhere.
Group coefficients of seeds having the same closure:
\[
 p(M)=\sum_{K\in\mathcal M}\gamma_K\,1[K\subseteq M].
 \tag{12}
\]
Since $p$ is one somewhere, at least one $\gamma_K$ is nonzero.
Choose an inclusion-minimal such $K$. Evaluating (12) at $K$ gives
$p(K)=\gamma_K=1$: any other contributing closure would have to be a
strict subset of $K$ with nonzero coefficient, which minimality excludes.
Some seed of size at most $d$ generates this $K$, since otherwise
$\gamma_K$ would be zero. This proves (11).
\end{proof}

\begin{corollary}[Uniform search bound]\label{cor:seedsearch}
If $m\le\poly(N)$ and least-extension computation and acceptance evaluation
take polynomial time in the explicit input size $N$, an accepted member, when one exists,
can be found in
\[
 \poly(N) B(m,d),\qquad
 B(m,d)=\sum_{i=0}^{\min(d,m)}\binom mi,
 \tag{13}
\]
using polynomial workspace.
\end{corollary}

\begin{proof}
Enumerate seeds by increasing size, compute each consistent closure
(including the empty set), and evaluate acceptance.
Lemma \ref{lem:seed} guarantees success by size $d$.
The algorithm does not need the expanded polynomial. For successful
witness search it does not even need $d$ as input: it continues until a
witness is found. To stop early and declare absence, it needs a justified
degree bound, or must enumerate all seeds. Seeds and closures can be
streamed, yielding polynomial workspace.
\end{proof}

For Horn CNF, models are intersection closed and forward chaining computes
the least extension or detects contradiction.
For a consistent rank-$d$ affine system, retain independent equations
$\ell_j(x)=b_j$. Its acceptance predicate is
\[
 a(x)=\prod_{j=1}^d(1+\ell_j(x)+b_j)
 \quad\text{over }\F,
 \tag{14}
\]
whose multilinear degree is at most $d$.
This gives (13) for Horn plus rank-$d$ affine constraints.
For satisfiable 2-CNF, first find a model $a$ and substitute $x=a\oplus y$.
The resulting 2-CNF is satisfied by zero, so every nontrivial clause
has a negative literal and is Horn. Shift the affine right-hand side
accordingly. Unsatisfiable 2-CNF is detected before the transformation.

These tractability regimes are established in the literature.
Brakensiek, Gopi, and Guruswami give an
$n^{k(M-1)+O(1)}$ algorithm for Horn with $k$ modular equations for
prime-power $M$, including $M=2$ \cite{brakensiek2019}.
Arvind and Guruswami study 2-CNF under affine-subspace constraints,
including its unrestricted and parameterized hardness
\cite{arvind2021}. These results do not give an unconditional
superpolynomial lower bound for our explicit family.

\begin{theorem}[Closure obstruction for a sequential witness chain]\label{thm:chainclosure}
Suppose one uniform algorithm reconstructs the chain from its canonical
input, and, for every prefix it actually selects, obtains an intersection-closed
extension representation with $m_j$ coordinates and an acceptance predicate
of degree at most $d_j$ on that family.
Assume every selected valid prefix extends, and all representation
acquisition, closure computations, predicate evaluations, witness lifts,
and the number of stages have one uniform polynomial bound in final
length $N$. Then
\[
 H(F_J)\le
 \left[\poly(N)\sum_{j<J}B(m_j,d_j)\right]^{D_0}
 \tag{15}
\]
for a fixed simulation exponent $D_0$.
In particular, if $m_j\le N^{O(1)}$ and $d_*=\max_{j<J}d_j$,
\[
 H(F_J)\le N^{O(d_*+1)}.
 \tag{16}
\]
\end{theorem}

\begin{proof}
Use Corollary \ref{cor:seedsearch} on the current actual prefix, lift the
accepted closure to a valid next witness, and retain only this witness.
Every selected prefix extends, so no backtracking is needed.
Sum the uniform costs, reconstruct and verify the final assignment,
and account for fixed polynomial simulation overhead. The same fixed
program performs every stage. This proves (15); the elementary bound
$B(m,d)\le(m+1)^{d+1}$ gives (16).
\end{proof}

A sharper statement retains the binomial count. Put
$\delta_j=\min(d_j,m_j)$, and interpret the expression below as zero
when $\delta_j=0$:
\[
 E_j=\min\left\{m_j,\ 
    \delta_j\log\left(\frac{e m_j}{\delta_j}\right)\right\}.
\]
The standard bounds $B(m,d)\le2^m$ and
$B(m,d)\le(em/\delta)^\delta$ yield
\[
 \log H(F_J)\le O\left(\log N+\max_{j<J} E_j\right).
 \tag{17}
\]
Therefore a chain satisfying (4) must have
\[
 \max_{j<J}E_j=\Omega((\log N)^2).
 \tag{18}
\]
In particular, $m_j=O(\log N)$ for every stage forces polynomial work,
whatever the degrees; and $d_*=o(\log N)$ with polynomially many
coordinates rules out (4). Thus $d_*=\Omega(\log N)$ is another
necessary condition within this represented extension architecture.

The degree here is the degree of the \emph{acceptance predicate on the
closure family}; a matrix rank is only a usable upper bound in (14).
Large rank alone need not mean large search cost.
The current project note's connected selector control already has
$3^m$ surviving closure-polynomial terms and rank $m$, yet an explicit
simple witness and treewidth at most three \cite{atc2026}.
An expensive expanded representation does not require an algorithm
to construct that representation.

\section{Implications for the polynomial and quasipolynomial hypotheses}

Theorems \ref{thm:collision} and \ref{thm:chainclosure} are upper bounds.
They are consistent with both $\mathsf P=\mathsf{NP}$ and
$\mathsf P\ne\mathsf{NP}$. They invalidate proposed hardness statements
for the covered construction regimes, not either complexity universe.

For a fixed desired lower degree $p>1$,
\[
 H(F_N)\ge2^{\Omega((\log N)^p)},
 \tag{19}
\]
the same comparison gives necessary conditions
\[
 R=\Omega((\log N)^p),\qquad
 \max_jE_j=\Omega((\log N)^p).
 \tag{20}
\]
When $m_j\le\poly(N)$, the second implies
$d_*=\Omega((\log N)^{p-1})$. The conditions apply only if the respective
collision or closure representation is available with the charged
hypotheses. A candidate need not have both representations.

If (19) were established for the actual benchmark, (2) would exclude
SAT algorithms of time $2^{O((\log N)^q)}$ for every $q<p$.
For $p=2$ this excludes all polynomial algorithms and permits consistency
with a sufficiently large quasipolynomial upper bound.
It does not exclude the whole union of quasipolynomial time classes.
No ETH assumption or claimed ETH refutation is used here.

\section{A second attempted contradiction: finding errors in short solvers}

In this section only, $N$ denotes the exact input length; replacing it
by $N+2$ elsewhere changes no asymptotic bound.
Fix an integer clock exponent $k$. Let $S_{p,N,k}$ simulate $p$ for
$N^k$ steps on an input of length $N$, treating timeout as failure.
For the explicit SAT verifier $V$, define
\[
 \operatorname{Err}_{N,k}(p,x,y)=
 V(x,y)\wedge\neg V(x,S_{p,N,k}(x)).
\]
The short-program error language is
\[
 L_k=\{1^N\#p:\ |p|\le b_0+\lfloor\log N\rfloor,\
       \exists x,y,\ |x|=N,\ \operatorname{Err}_{N,k}(p,x,y)\}.
 \tag{21}
\]
For every fixed $k$, $L_k$ is in NP and is sparse.
The verifier has fixed polynomial runtime and there are only
polynomially many admitted codes through each input length.
By Mahaney's theorem, a polynomial many-one reduction from SAT to $L_k$
would imply $\mathsf P=\mathsf{NP}$ \cite{mahaney1980}.
Trying to hardcode an arbitrary $n$-bit formula into a program with
$O(\log N)$ bits generally forces exponential $N$, so the unary output
ceases to be a polynomial-size reduction.

A polynomial-size nonuniform table of counterexamples exists:
for every erroneous admitted code select one pair $(x,y)$.
This counts only representation size. It supplies neither a polynomial
uniform acquisition algorithm nor short proofs that omitted programs
are correct. For a correct program, the corresponding witnessing
implication still expresses global clocked correctness at that length.
This is the earlier reflection gap, not a solution to it.
Also, $L_k\in\mathsf P$ by itself does not give a generic prefix-search
self-reduction confined to the special short-description format.

\begin{proposition}[The refuter must pay for the candidate clock]\label{prop:refuter}
Use an ordinary effective programming language with a fixed-point
construction having fixed polynomial simulation overhead.
Assume there are at least two distinct satisfiable encodings at
arbitrarily large admissible lengths $N$.
There is no single total deterministic refuter
$R(1^N,1^k,p)$ that, for every $k,N$ and admitted $p$, returns a valid
same-length pair witnessing (21) whenever such a pair exists, and
whose runtime is bounded by one polynomial in $N+k+|p|$ with degree
independent of $k$.
\end{proposition}

\begin{proof}
Suppose such an $R$ exists. By a standard self-quoting fixed-point
construction, define a fixed program $p_{R,k}$ that, on input $x$
of length $N$, runs $R(1^N,1^k,p_{R,k})$. If the returned pair is
$(x^*,y^*)$ and $x=x^*$, return $y^*$; otherwise fail.
The polynomial degree of this program, including fixed simulation
overhead, is independent of the fixed constant $k$.
Choose $k$ larger than that degree. For sufficiently large admissible
$N$, this fixed code fits the logarithmic description budget and
finishes within the candidate clock $N^k$.

The program can succeed on at most one input at that length.
Because there are two distinct satisfiable encodings, it has an error.
The refuter must therefore return a valid error pair $(x^*,y^*)$.
But on input $x^*$ the program returns exactly $y^*$ within its clock,
so that pair is not an error. This contradiction proves the proposition.
\end{proof}

The proposition does not prohibit
$\forall k\,\exists R_k$ with exponent depending on $k$.
Nor does it prohibit a procedure polynomial in the unary clock
$N+N^k$. Under a hypothetical SAT algorithm of time $N^a$,
compiling and searching the error predicate can cost $N^{O(ka)}$.
It does not fit the candidate's $N^k$ clock in the way the proof
requires. Thus the self-reference creates no contradiction with
$\mathsf P=\mathsf{NP}$.

The deterministic dreambreaker theorem of Bogdanov, Talwar, and Wan
assumes $\mathsf P\ne\mathsf{NP}$ and guarantees certified errors
infinitely often with resources charged to the candidate running time
\cite{bogdanov2010}. It cannot be inserted as an unconditional lemma
inside an argument starting from $\mathsf P=\mathsf{NP}$.

\section{Computational verification and provenance}

Both programs are deterministic and use exact Boolean arithmetic.
No randomized seed is used. Each invocation completed well below
the repository's 30-second budget.
\begin{description}
\item[] Producer:
\path{experiments/review/witness_amplifier_controls_20261111.py}.
SHA-256:
\texttt{9197bc468d255e7746f6d2f6c8a421c6186bce887584ef12212c8cd26c6a5ff8}.
Runtime in the recorded run: approximately $0.245$ seconds.
\item[] Independent audit:
\path{experiments/review/witness_amplifier_independent_audit_20261111.py}.
SHA-256:
\texttt{55a09519a0883617bd891b1d65cbf5ace4c3e24b33cd6723b49077d1a69edbdc}.
The audit does not import the producer.
Runtime: approximately $0.0043$ seconds.
\item[] Raw directory:
\path{results/raw/review/2026-11-11_WITNESS_AMPLIFIER_CONTROLS/}.
Files: \path{witness_amplifier_controls.json},
\path{independent_audit.json}, \path{nested_feistel_h2_k2.cnf},
and \path{nested_feistel_h2_k2_witness.json}.
\end{description}

\begin{center}
\begin{tabular}{rrrrr}
\toprule
Half-width $h$ & Depth $k$ & Variables & Clauses & Inverse evaluations\\
\midrule
2&1&104&406&2\\
2&2&201&797&4\\
3&2&299&1193&4\\
4&3&592&2376&6\\
5&4&983&3955&8\\
\bottomrule
\end{tabular}
\end{center}

The producer checks 4,080 permutation/inverse pairs and 24 elementary
gate truth cases. It exhausts 2,048 primary assignments on the smallest
nonlinear chain and finds exactly 16 satisfying primary assignments,
with no disagreement between the direct collision semantics and the
CNF construction. This is not enumeration of all assignments to all
104 variables: exact gate equivalences justify completing auxiliaries.
The independent audit additionally checks 96 proper gate/padding cases,
48 aliased-input cases, all four padding settings of the exported
201-variable/797-clause witness, and 64 inputs across all four two-bit
keys of a separate four-bit permutation control. All pass.
The last control has coordinate algebraic degrees two or three,
confirming nonlinearity.

For the closure lemma, the producer enumerates all 121 nonempty
intersection-closed families on three coordinates and all polynomial
coefficient choices of degree at most $0,1,2,3$.
Across 48,642 family/polynomial pairs it finds zero mismatches between
existence of an accepted model and existence of an accepted
small-seed closure. This is finite evidence supporting the separately
proved lemma, not an asymptotic proof or a proof-assistant check.
Neither program measures minimum $H$ or establishes a lower bound.

Independent temporary-agent review checked the collision theorem and
canonical-recognition hypothesis, the closure proof and sharper binomial
bound, and the self-repair clock quantifiers. The producer and independent
CNF auditor are separate implementations. This is mathematical review
and executable finite validation, not formal proof-assistant verification.

\section*{CORPUS REVIEWED}
The run reviewed the current root manuscript, AGENTS.md, CODEX.md,
the repository research skill and mode/protocol references, both
root proof trackers, the November 9 governing coordination handoff,
targeted MAIN and specialist history, and material exploration/source
artifacts. Root manuscript was pinned to blob
\texttt{ca1c40a11155442a7ecdf74d0e05db3f83abae9a}.
This is a targeted corpus reconciliation, not a claim to have reread
every historical byte or run tests on a full repository checkout.

Material additions to the previous run's source set:
\begin{description}
\item[] The actual October 13 uploaded minimum-work note,
SHA-256
\texttt{3dd480be5ed2450447519565ba2a23f596f4cf19cc8ce78e47425aeff63e654d}.
\item[] AFFINE-TABLE-COMPOSITION2,
\path{research/explorations/sat-composition/2026-11-11_AFFINE_TABLE_COMPOSITION2/compatibility_continuation.tex},
blob \texttt{3809296b0a5f6ea85fc03bf1d8441b04e3017585}.
\item[] LSPV-1/586,
\path{research/agents/lspv-1/2026-11-11_1857_LSPV1_CHECKER_ACCEPTANCE_CEGAR_NOGO586.md},
blob \texttt{129ae6001f1067244458f1dc5f58f705441216ae}.
Its reported accepted-byte semantic counterexample is not independently
rerun in this continuation.
\item[] The uploaded block-protection/GC draft and Type-I pruning notes.
Their structural or conditional algorithm statements do not supply
the all-program lower bound in (3).
\end{description}

\section*{DUPLICATION / CONTRADICTION AUDIT}
The October 13 note already gives a fixed-seed replay obstruction and
the constant-output collision-search objection. The new contribution
is the adaptive all-prefix accounting, explicit growing-width thresholds,
the implemented nonlinear inverse control, and the closure-degree
transfer to the same minimum-work target.
Horn/affine tractability and the closure basis are prior results;
this report does not rename them as a new complexity separation.

The counterexample table, sparse-language consequence, and clock
self-repair proposition block stronger witnessing claims. None
contradicts a fixed-degree polynomial SAT assumption.

\section*{WHAT CHANGED RELATIVE TO ROOT MAIN}
The quantitative exclusion region is enlarged.
The root's OPEN all-solver arrows remain OPEN.
The latest /341 fixture acceptance defect strengthens the need for
source and inference soundness before any proof-system transfer.
It does not affect the independently verified positive-witness
benchmark (1). No root ownership region is modified by this Codex run;
the handoff is for the synthesis owner.
\textbf{TRACKERS REVIEWED; NO MATERIAL CHANGE} to a load-bearing status.

\section*{ADOPTED / REJECTED / RETIRED PRIOR COMPONENTS}
\begin{description}
\item[] Adopted: the $H$ bridge, fixed-seed replay warning, exact circuit
encoding, Horn closure, affine-rank reduction, and paid prefix interfaces.
\item[] Refined: constant-output enumeration to adaptive chains and
output-width thresholds; affine rank to degree on a closure family;
coarse seed counts to the binomial bound (17).
\item[] Falsified for the explicit candidate: a growing minimum-work
multiplier produced by truncated efficiently invertible Feistel layers.
\item[] Rejected: multiplying stage costs merely because queries depend
on earlier answers; inferring work from a large expanded polynomial;
using a black-box separation as an explicit-machine lower bound;
using one clock-independent refuter without paying its time.
\item[] Stale or overstated: any use of the archived /341 checker's
raw ACCEPT as universal semantic soundness. The source report already
identified that defect; this continuation does not claim a new checker
counterexample.
\end{description}

\section*{Hostile controls and barrier audit}
\begin{description}
\item[] Horn, 2-SAT and affine: exact algorithms give the closure bound.
High affine rank alone is harmless for Gaussian elimination.
\item[] Nonlinearity and sequential dependence: the Feistel control
contains both and still has polynomial witness construction.
\item[] Padding and auxiliary variables: all short clauses are expanded
exactly; independent controls check padding and aliased gates.
\item[] Representation and wrappers: one fixed program reconstructs
the canonical family; arbitrary descriptions or a free decomposition
are not assumed. Some-prefix existence is not all-prefix extension.
\item[] Relativization: the collision enumeration upper bound uses only
evaluation and survives adding an oracle with its actual query costs
charged. The cited black-box lower bound has a different model.
\item[] Algebrization: the closure argument is finite Boolean algebra.
Extra algebraic access cannot invalidate an explicit upper algorithm;
it gives no separation against such access.
\item[] Natural Proofs: no constructive property of arbitrary Boolean
circuits is used to claim a general circuit lower bound. No barrier
avoidance is asserted.
\item[] Cryptographic transfer: efficient Feistel inversion is an exact
algorithmic transfer. Hardness of a hash or permutation is not assumed.
\end{description}

\section*{Status table}
\begin{longtable}{p{0.54\textwidth}p{0.20\textwidth}p{0.18\textwidth}}
\toprule
Claim & Status & Scope\\
\midrule
Adaptive collision upper bound (6) & PROVED/EXACT & Paid explicit family\\
Feistel polynomial witness algorithm (8) & PROVED/EXACT & Explicit construction\\
Closure small-seed lemma (11) & PROVED/EXACT & Intersection-closed family\\
Closure-chain bound (15)--(17) & PROVED/EXACT & Actual-prefix extensions\\
Short-program language sparsity & PROVED/EXACT & Fixed clock exponent\\
Clock-independent refuter impossibility & PROVED/EXACT & Proposition \ref{prop:refuter}\\
Finite CNF and closure tests & COMPUTATIONAL & Recorded finite domains\\
Growing work multiplier (3) & OPEN & Every admitted witness program\\
P versus NP & OPEN & Unrestricted deterministic time\\
\bottomrule
\end{longtable}

\section*{Single next target}
For the next \emph{fully specified} candidate-dependent wrapper,
construct its exact witness-extension relation and determine whether
the actual selected prefixes are all extendible.
Then apply (5) and (15) to that relation with representation acquisition
and witness lifting charged. A surviving candidate must either defeat
these algorithms by a proved parameter bound, or exhibit a genuine
compatibility obstruction requiring revisiting earlier witnesses.
Neither condition alone supplies (3). The missing theorem remains an
all-program lower bound for the same explicitly defined $H$.

\bibliographystyle{plain}
\bibliography{research/agents/codex/2026-11-11_2017_ADAPTIVE_COLLISION_AND_CLOSURE_BARRIERS_refs}
% END_OF_VERIFIED_REPORT::CODEX::ADAPTIVE-COLLISION-CLOSURE-BARRIERS
\end{document}
